\chapter
[\CO{[NEU]} Neurologische Störungen]
{\CC{[NEU]} Neurologische Störungen}
\ChapterIntro
{%
Neurologische Störungen sind Störungen bzw.
Erkrankungen, Notfälle und sonstige 
die das Nervensystem betreffen. Sie
gehen oft mit einer Störung des Bewusstseins oder anderer
Funktionen des Nervensystems einher. Die \EE{Ursachen} von
neurologischen Erkrankungen sind sehr vielfältig. }
{%
\begin{description}
  \item[Maintainer:] Sebastian Gabriel
  \item[Autoren:] Diverse
  \item[Reviewer:] Standard-Reviewprozess
  \item[Version:] {\DocVersion} ({\DocVersionInfo})
  \item[SHA1:] {\DocGitCommit}
\end{description}

{\TxTeilkapitelAASS}
}

\bigskip

\noindent Die folgende Übersicht nennt einige
Beispiele für Ursachen für neurologische Störungen, ohne
Anspruch auf Vollständigkeit:

\begin{WidePar}
\begin{multicols}{2}
\begin{itemize}
  \item \E{Schlaganfall}
  \item \E{Vergiftungen} können fast alles verursachen, inkl.
  neurologischer
  Symptome.
  \item \E{Schädel-Hirn-Trauma} (SHT); Ein Schlag auf
  den Kopf erhöht \emph{nicht} das Denkvermögen \ldots
  \item \E{Stoffwechselstörungen} welche Einfluss auf das
  Nervensystem nehmen, z.\,B. Zuckerentgleisungen
  \item \E{Infektionen} können das Nervensystem beeinflussen,
  z.\,B. Verwirrtheit bei Fieber
  \item \E{Alter} beeinflusst auch das Nervensystem, im Rahmen
  einer Demenz kommt es z.\,B. zum Abbau der Denkleistung
  \item \E{Krampfanfälle} können vom zentralen Nervensystem ausgehen
  \item \E{Erkrankungen des Stützapparates} Aufgrund der
  räumlichen Nähe zum Stütz- und Bewegungsapparates kann das
  Nervensystem in Mitleidenschaft gezogen werden, z.\,B. bei
  Bandscheibenvorfällen
\end{itemize}
\end{multicols}
\end{WidePar}

\clearpage
\Layout{2}{{\TxQuerverweise}}{%
\begin{itemize}
  \item Infektiöse (bakterielle) Meningitis: \dotfill
  \FullPrettyRef{topic:meningitis}
\end{itemize}
}{}{}{}{}




\section{Schlaganfall, Apoplektischer Insult, Hirninfarkt}
\index{Schlaganfall}
\index{Insult}
\index{Apoplexie}
\index{Apoplektischer Insult}
\label{topic:tia}
\label{topic:insult}
\label{topic:schlaganfall}

\HeadingSub{\Lang{Syn.} Apoplektischer Insult, Apoplexie}

% \Definition{Siehe Arten.}

\Layout{0}{Arten}
{%
Grundsätzlich gibt es 2 Formen, die aber vor Ort nicht
unterschieden werden können:
\begin{itemize}
  \item \EE{\T{Trockener Insult}}:
  \I[Ischämie!trockener Insult]{Ischämie}. Verschluss eines
  Hirngefäßes mit anschließendem Absterben der versorgten Hirnregion
  \item \EE{\T{Blutiger Insult}}: \I{Hirnblutung}, z.\,B.
  durch Platzen eines Aneurysmas eines Hirngefäßes, kann spontan
  oder aufgrund eines Traumas auftreten.
\end{itemize}

}
{
\begin{itemize}
  \item \enquote{Trockener} Insult
  \item \enquote{Blutiger} Insult
\end{itemize}

}{}{}{}

\Layout{0}{\T{TIA}}
{%
\index{TIA}
\index{Transitorisch-Ischämische Attacke}
\label{topic:tia}
TIA ist die Abkürzung für Transitorisch-Ischämische Attacke.
Von einer TIA spricht man, wenn die Symptome  innerhalb von
24 Stunden vergehen. Solange noch Symptome bestehen, kann
man nicht von einer TIA sprechen, der Patient ist wie bei
jedem anderen Schlaganfall zu behandeln.
\bigskip
}
{
% TODO MED: Slide fehlt!
\begin{itemize}
  \item Transitorisch-Ischämische Attacke.
  \item Symptome <\,24\Hour*.
\end{itemize}

}{}{}{}

% \Dias
% {
%     \includegraphics[width=\linewidth]{./images/noauth/AASdev20090228-img83.png}
%     \Captionof{figure}{\AuthNotesPicLegalNotOk{img83}{}Insult: Eine spastische
%     Halbseitenlähmung kann sich als Spätfolge Wochen nach einem durchgemachten
%     Schlaganfall e ntwickeln.
%     \SourcePicture{Quelle nicht eruierbar.}}
% }{
%     \includegraphics[width=\linewidth]{./images/noauth/AASdev20090228-img84.png}
%     \Captionof{figure}{\AuthNotesPicLegalNotOk{img84}{}Insult:
%     Halbseitenlähmung im Gesicht.
%     \SourcePicture{Quelle nicht eruierbar.}}
% }{}


\Layout{0}{\TxSymptome}
{%
\label{topic:hirndruckzeichen}
Die diversen \EE{Funktionsareale}
(Sprachzentren, Sehzentrum, motorische Zentren, \ldots)
sind in verschiedneen Regionen des Hirns verteilt. Je
nachdem welches Gebiet durch die Ischämie oder Blutung
betroffen ist, können die Symptome sehr unterschiedlich
sein.  Typische Störungen sind:

\begin{itemize}
  \item \EE{Motorik}: Störungen der Kraft und Beweglichkeit
  sind sehr häufig.
  Charakteristisch für einen
  Schlaganfall sind sog. \T{Halbseitenzeichen}: Es kommt
  dabei zu einer \E{einseitigen Kraftminderung}, dies kann
  bis zu einer \E{schlaffen, kompletten Lähmung} führen. Ein
  \E{herabhängender Mundwinkel} ist ein Zeichen einer
  Lähmung der
  Gesichtsmuskulatur. Durch die Lähmungen kann es zu
  \EE{Schluckstörungen} kommen, es besteht dann eine erhöhte
  Aspirationsgefahr.
  
  Der \T{Herdblick} ist eine Blickdiagnose: Es kann zu
  Blicklähmungen kommen, der Patient blickt in Richtung des
  \enquote*{Herdes des Geschehens}.
  
  \item \EE{Sprache}: Störungen der Sprache sind ebenfalls
  sehr häufig. Oft nimmt man eine \EE{verwaschene Sprache} wahr,
  welche (fälschlich) an Trunkenheit erinnert. \Z{Weiters
  kann es zu einer \enquote{wirren Sprache} kommen: Worte fehlen dabei
  oder Silben werden vertauscht. Im Extremfall kann die
  Sprache gänzlich \enquote{verloren} gehen. Auf der anderen Seite
  kann die Sprache zwar erhalten, aber das
  Sprach\E{verständnis} geschädigt sein.}
  
  \item \EE{Wahrnehmung}:
  Es kann zu Wahrnehmungstörungen wie z.\,B.
  \EE{Sehstörungen}, plötzliche Erblindung, Doppelbilder,
  Gesichtsfeldausfälle kommen.
  
  \item \EE{Bewusstsein}: Störungen des Bewusstseins sind
  gefährlich und müssen bei Insult-Patienten immer erwartet
  werden. Der Bewusstseinszustand kann sich auch plötzlich
  verschlechtern.
  
  \item \EE{Allgemein}: Weiters können unspezifische
  Symptome wie z.\,B. Schwindel und Kopfschmerzen auftreten. Ein
  plötzlicher, \EE{donnerschlagartiger Kopfschmerz} ist ein
  klassisches Symptom für eine spontane Hirnblutung.
  
  \item \EE{\T{Hirndruckzeichen}}: Bei einer Erhöhung des
  Hirndrucks, z.\,B. aufgrund einer Hirnblutung, kommt es zu
  einer typischen Kombination von Symptomen: 
  
  \parpic[4cm][r]{\includegraphics[width=4cm]{./images/trauma/pupillendifferenz.jpg}}
  Meist das erste Symptom sind \EE{Kopfschmerzen}, gefolgt
  von Übelkeit. In weiterer Folge treten
  \EE{Bewusstseinsstörungen} auf. Der Patient ist
  \EE{hyperton, aber bradykard} (\EE{RR\,\Sign{↑},
  HF\,\Sign{↓}}), der Puls ist langsam, aber äußerst stark
  spürbar (\I{Druckpuls}). Weiters ist die \EE{Ungleichheit
  der Pupillen} \Z{(Anisokorie)\index{Anisokorie|Z}}, bzw.
  eine verlangsamte Lichtreaktion
  typisch.\ZFootnote{\E{Erhöhter Hirndruck}:In schweren
  Fällen kommt es zu sog. Strecksynergismen, d.\,h.  der
  Patient reagiert bei Schmerzreiz nur mit einem
  ungerichtetem Strecken der Extremitäten.}
  \nocite{StriebelAnaesthesie2} %[Bd.2, S.1344]
  
  \Z{Hirndruckzeichen werden nicht nur durch
  Hirnblutungen verursacht, auch z.\,B. beim \Z{Hirnödem}
  oder Schädel-Hirn-Trauma (\Abk{SHT}) können sie
  vorliegen.}
\end{itemize}

Präklinisch kann man nicht sicher zwischen einer Ischämie
(trockener Insult) und einer Blutung unterscheiden.
Auf jeden Fall ist die \EE{Bewusstseinslage} sehr engmaschig
zu überwachen; eine rasche Verschlechterung ist jederzeit
möglich! 
% Bei einer \E{Blutung} kommt es jedoch
% oft zu einer \E{raschen Verschlechterung des Zustandes}, 
% insbesonders hinsichtlich des Bewusstseinsgrades.

%\begin{multicols}{2}
% \begin{labeling}[~--]{xxxxxxxxxxxxxxxxxx}
%   \item[Motorik] Störungen der Kraft und Beweglichkeit sind sehr häufig:
%   
%   \T{Halbseitenzeichen}: Aufgrund des Ausfalls eines Zentrums im
%   Gehirn kann es zu einer \E{Kraftminderung} einer Seite kommen, dies kann bis
%   zu einer \E{schlaffen, kompletten Lähmung} führen.
%   
%   Ein \E{herabhängender Mundwinkel} ist ein Zeichen einer Lähmung
%   der Gesichtsmuskulatur.
%   
%   Durch die Lähmungen kann es zu \EE{Schluckstörungen} kommen, es
%   besteht dann eine erhöhte Aspirationsgefahr.
%   
%   Der \T{Herdblick} ist eine \enquote{Blickdiagnose}: Es kann zu Blicklähmungen
%   kommen, der Patient blickt in Richtung des \enquote{Herd des Geschehens}.
%   
%   \item[Sprache] Störungen der Sprache sind ebenfalls sehr
%   häufig
%   
%   Oft nimmt man eine \EE{verwaschene Sprache} wahr, welche (fälschlich) an
%   Trunkenheit erinnert.
%   
%   Weiters kann es zu einer \enquote{wirren Sprache} kommen: Worte fehlen dabei oder
%   Silben werden vertauscht. Im Extremfall kann die Sprache gänzlich \enquote{verloren}
%   gehen. Auf der anderen Seite kann die Sprache zwar erhalten, aber das
%   Sprach\E{verständnis} geschädigt sein
%   
%   \item[Wahrnehmung]
%   Es kann zu Wahrnehmungstörungen wie z.\,B. \EE{Sehstörungen}, plötzliche
%   Erblindung, Doppelbilder, Gesichtsfeldausfälle kommen.
%   
%   \item[Bewusstsein] Störungen des Bewusstseins sind gefährlich und müssen bei
%   Insult-Patienten immer erwartet werden. Der Bewusstseinszustand kann sich
%   auch plötzlich verschlechtern.
%   
%   \item[Allgemein] Weiters können unspezifische Symptome wie z.\,B.
%   Schwindel und Kopfschmerzen auftreten.
%   Ein plötzlicher, \EE{donnerschlagartiger Kopfschmerz} ist ein klassisches
%   Symptom für eine spontane Hirnblutung.
%   
% \end{labeling}
%\end{multicols}
}
{
Je nach betroffener Hirnregion sehr unterschiedlich.
\begin{itemize}
  \item Halbseitenzeichen (Schwäche, Lähmung)
  \item Herdblick
  \item Verwaschene Sprache
  \item Sehstörungen
  \item Bewusstseinsstörungen
  \item Kopfschmerzen
  \item Hirndruckzeichen
  \begin{itemize}
    \item Bewusstseinsstörungen
    \item RR\,\Sign{↑} \& HF\,\Sign{↓}
    \item Ungleichweite Pupillen, verlangsamte Lichtreaktion
    \item Strecksynergismen
  \end{itemize}
\end{itemize}
}{}{}{}


\Layout{0}{Differentialdiagnose}
{
Die wichtigste Differentialdiagnose ist eine Störung des
\EE{Blutzucker}-Stoffwechsels. Sowohl die Hyper-, als auch
die Hypoglykämie können neurologische Symptome haben, daher
muss der Blutzucker immer gemessen werden! Ebenso ist immer
an \EE{Vergiftungen} zu denken.

Die Abgrenzung zu anderen neurologischen Erkrankungen ist
ohne weitergehende Untersuchungen oft schwierig bzw. schwer
möglich. 
}{
\begin{itemize}
  \item Blutzucker ${\downarrow}$, ${\uparrow}$
  \item Sonstige neurologische Erkrankung.
  \item Vergiftungen
\end{itemize}

}{}{}{}

\MassnahmenMK{Spezielle
Maßnahmen}{Insult}{m:insult}
{\EE{Taktik}: \Speech{\Ca{Vitale Bedrohung!}
Zeitkritisch (>>\EE{Time is brain.}<<):
Vitalfunktionen sichern, rascher Transport an eine geeignete Abteilung.}

Achtung: Ausnahme bei der Notarzt-Nachforderung
}{}
% \MassnahmenNG{Spezielle
% Maßnahmen}{Insult}{\label{m:insult}}
% {\EE{Taktik}: \Speech{\Ca{Vitale Bedrohung!}
% Zeitkritisch (>>\EE{Time is brain.}<<):
% Vitalfunktionen sichern, rascher Transport an eine geeignete Abteilung.}
% 
% \begin{itemize}
%   \item \TxMassVitBes
%   \begin{itemize}
%     \item Lagerung: \E{leicht erhöhter Oberkörper}
% %     \item \ce{O2}-Gabe gemäß \prettyref{m:sauerstoffberieselung}
%     
%     \item Notarzt: Grundsätzlich ist ein
%     Schlaganfall-Patient aufgrund seiner Diagnose als vital
%     bedroht einzustufen. Es gibt jedoch eine Besonderheit:
%     Eine definitive Therapie in einer Spezialabteilung
%     (Stroke Unit) kann andere, noch nicht abgestorbene Teile
%     des Hirns vor weiterem Schaden bewahren, d.\,h. es ist
%     sehr wichtig, dass diese Therapie möglichst schnell
%     erfolgt.
%     
%     Daher soll von einer
%     \EE{Nachberufung eines Notarztes \underline{abgesehen} werden,
%     \E{wenn folgende Umstände gegeben sind}}:
%     \begin{itemize}
%       \item Die Verdachtsdiagnose gilt als \E{sehr sicher}.
%       \item Der Patient ist sonst als \E{stabil} eingestuft
%       (\prettyref{m:ersteinschaetzung}) und
%       bewusstseinsklar\footnote{Bei der Einschätzung des Bewusstseins
%       darf man sich nicht durch die durch den Schlaganfall
%       verursachten Symptome (Sprachstörungen, Blickstarre \ldots)
%       beirren lassen.}.
%       \item Es gibt sonst \E{keine anderen Umstände} die die Berufung
%       eines Notarztes erforderlich machen.
%       \item Es muss versucht werden, eine Aufnahme auf eine
%       \E{Spezialabteilung} für Schlaganfälle zu
%       organisieren\footnote{Der Erfolg ist abhängig von der
%       Verfügbarkeit eines entsprechenden freien Bettes.} (Stroke Unit,
%       Kontaktaufnahme mit der Leitstelle).
%     \end{itemize}
%     \A{HART: Muss es unbedingt eine Spezialabteilung sein?
%     \enquote{bewusstseinsklar} könnte missverstanden werden.}
%     \A{\par Diskussion: }
%     \ACh{KOCH}{2010-02-13}{ok.}
%     \ACh{HART}{2010-02-13}{Muss es unbedingt eine Spezialabteilung sein?
%     \enquote{bewusstseinsklar} könnte missverstanden werden.}
%     \ACh{GABS}{2010-02-13}{Fußnote wegen Bewusstseinseinschätzung
%     eingefügt. Spezialabteilung: muss nur \underline{versucht} werden
%     eine Stroke Unit zu organisieren, je nach Verfügbarkeit.}
%     \ACh{GABS/STEU}{2010-04-06}{Ticket 38:
%     Besprechung gestern mit STEU: kann so übernommen werden, sieht
%     allerdings Problem wegen Inkonsistenz und für den Nutzer
%     möglicherweise verwirrender wenn-dann-Verzweigung. Bei Hirnblutung
%     ist das Vorgehen ein Nachteil. (Einwand GABS: Ischämisch häufiger,
%     Abwägung Benefit Narkose, etc. vs. schnelles CT)}
%   \end{itemize}
%   \item \EE{Neurocheck} inkl. Messung des \EE{Blutzuckers}!
%   \item Transportentscheidung: Wenn verfügbar: \EE{Stroke Unit}, sonst
%   Abklärung mit der Leitstelle wegen Alternativen (Abt. für
%   Neurologie oder Innere Medizin, mit Möglichkeit zur
%   Intensivüberwachung).
%   \item Bei Anforderung einer Stroke Unit ist anzugeben:
%   \begin{itemize}
%     \item Alter
%     \item Beginn der Symptome
%     \item Vorherige Schlaganfälle
%    \end{itemize} 
% \end{itemize}
% }

\section{Vom Gehirn ausgehende Krämpfe -- Zerebrale Krampfanfälle}
\label{topic:zerebrale-krampfanfaelle}
\label{topic:krampfanfall}

\Definition{Vom Gehirn ausgehende motorische Krampfanfälle aufgrund von akuten
oder chronischen Erkrankungen des Gehirns. Dabei entladen sich unkontrolliert
Nervenzell-Gruppen im Gehirn.}


\Layout{0}{Einteilung}
{%
Die zerebralen Krampfanfälle kann man auf verschiedene Arten
einteilen:

\begin{enumerate}
  \item Nach Art:
  \begin{itemize}
    \item \EE{Tonisch}: \enquote{Verkrampfung}, erhöhter
    Muskeltonus
    \item \EE{Klonisch}: Zuckungen
    \item \EE{Tonisch-klonisch}: Kombination aus beidem, das
    klassische Erscheinungsbild eines Krampfanfalles
  \end{itemize}
  \item Nach Lokalisation:
  \begin{itemize}
    \item \EE{Fokaler Anfall}: Der Krampfanfall ist auf
    eine Körperregion beschränkt, siehe
    \ref{topic:krampfanfall-fokaler}
    \item \EE{Generalisierter Anfall}: Er betrifft den
    ganzen Körper, siehe
    \ref{topic:krampfanfall-generalisierter}
    
  \end{itemize}
\end{enumerate}
}{%
\begin{itemize}
  \item Nach Art:
  \begin{description}
    \item Tonisch
    \item Klonisch
    \item Tonisch-klonisch
  \end{description}
  \item Nach Lokalisation:
  \begin{description}
    \item Fokaler Anfall
    \item Generalisierter Anfall
  \end{description}
\end{itemize}
}{}{}{}

\Layout{2}{Ursachen}
{%
Für zerebrale Anfälle gibt es eine Reihe von
unterschiedlichen Ursachen die direkt oder indirekt das
Gehirn betreffen. Einige betreffen die Hirn- und
Nervenstruktur, andere sorgen dafür, dass das Hirn zu wenig
Sauerstoff und andere Nährstoffe erhält. Die häufigsten
Ursachen und Grunderkrankungen sind im Folgenden aufgeführt:

\begin{itemize}
  \item \EE{Epilepsie}: Die Epilepsie ist eine chronische
  Erkrankung welche durch wiederkehrende Episoden von zerebralen Krampfanfällen gekennzeichnet ist.
  \item \EE{Vergiftungen}: (Medikamente/Drogen/...) können
  wieder fast alles machen, also auch Krampfanfälle. \citep{Supady2009}
  \item \EE{Hirntumor}: Durch den Tumor wird das Hirn
  beschädigt, auch dadurch kann es zu Anfällen kommen.
  \item \EE{Narbengewebe}: nach Insult oder Hirnverletzung.
  \item \EE{Hirnblutung}
  \item \EE{Sauerstoffmangel}: z.\,B. Kollaps, Ischämie
  aufgrund eines Gefäßverschlusses
  \item \EE{Zuckerstoffwechselstörung}: Der Dauerbrenner.
  Eine Blutzuckerentgleisung kann fast alle neurologischen Symptome vortäuschen!
  \item \EE{angeboren}: ohne organisch fassbare  Ursache
  \item \EE{Fieberkrampf}: Fieber bei Kleinkindern
  \item \EE{Infektionen}:
  \item \EE{Entzug}: Entzug von Alkohol oder anderen Drogen
\end{itemize}
}{
}{}{}{}






\subsection{Fokaler Anfall}
\index{Krampfanfall!fokaler}
\label{topic:krampfanfall-fokaler}
Der Krampfanfall ist auf eine Körperregion beschränkt,
z.\,B. auf eine Hand. Der Patient kann bei Bewusstsein sein.
Ein fokaler Anfall kann sich ausweiten und in einen
generalisierten Anfall übergehen.

\Z{Der fokale Anfall wird
auch \emph{\enquote{petit mal}}\ZFootnote{Petit mal: \emph{franz.:} \enquote{kleines
Unglück}} genannt.}
      


\subsection{Generalisierter Anfall}
\index{Krampfanfall!generalisierter}
\label{topic:krampfanfall-generalisierter}

\Layout{0}{Beginn}
{%
Er betrifft den ganzen Körper. Dem Anfall kann ein
\E{Dämmerzustand} 
vorausgehen, \Z{eine sog. \enquote{Absence}}. Der Patient wirkt
dabei plötzlich abwesend und \emph{\enquote{irgendwie komisch}}.
Dann verliert er das Bewusstsein und die Muskeln werden
anfangs schlaff. Wenn der Patient gestanden ist, fällt er
ungebremst auf den Boden (Cave: \EE{Verletzungen!}).
Manchmal kommt es vor, dass der Patient einen
\I{Initialschrei} von sich gibt, der
gruselig klingen kann. 
}{
\begin{itemize}
  \item Dämmerzustand -- Absence
  \item Initialschrei
  \item Evtl. Sturz
  \item Evtl. Folgeverletzungen
\end{itemize}

}{}{}{}

\Layout{0}{Krampf}
{%
Dann kommt es zu
\EE{tonisch-klonischen Krämpfen},
die den ganzen Körper betreffen. Während des Anfalls ist der
Patient nicht bei Bewusstsein, es besteht außerdem i.\,d.\,R.
ein Atemstillstand. Er kann jedoch starke Kräfte entwickeln, es
besteht massive
Verletzungsgefahr.
}{
\begin{itemize}
  \item Tonisch-klonische Krämpfe
  \item Verletzungsgefahr
\end{itemize}
}{}{}{}


\Layout{0}{Danach}
{%
Der Krampf dauert meist wenige Minuten und hört von selbst
auf. Der Patient ist danach erstmal bewusstlos, im weiteren
Verlauf mit mehr oder weniger großem Aufwand erweckbar. Er
befindet sich dann in der sogenannten
\T{Nachschlafphase}\index{Nachschlafphase}.

Der Patient \E{klart im Verlauf mehr und mehr auf}. Es ist
wichtig zu versuchen mit dem Patienten in Kontakt zu treten
um seinen Wachheitsgrad und dessen
\EE{Verlauf} beurteilen zu können. Für den Anfall selbst besteht
Amnesie! Daher muss man den Patienten aufklären was passiert
ist, wo er sich befindet, etc. Ein \EE{Neuro-} und
\EE{Traumacheck} ist nach \underline{\E{jedem}} Anfall
unbedingt notwendig! 
}{
\begin{itemize}
  \item Nachschlafphase
  \item \EE{Aufklaren}
  \item Amnesie
  \item Neurocheck
  \item \EE{Traumacheck} nach \underline{\E{jedem}} Krampfanfall
\end{itemize}
}{}{}{}


\Layout{0}{Sonstige Besonderheiten}
{%
Häufig kommt es während des Anfalles zu einem Zungenbiss,
außerdem zum Einnässen oder Einkoten.
}{
\begin{compactitem}
  \item Zungenbiss
  \item Einnässen, Einkoten
\end{compactitem}
}{}{}{}


\Z{Der generalisierte Anfall wird auch \emph{\enquote{grand
mal}}\ZFootnote{Grand mal: \emph{franz.:} \enquote{großes Unglück}} genannt.}



\Layout{0}{Status Epilepticus}
{\index{Status!epilepticus}
Der Status Epilepticus ist eine
lebensgefährliche Steigerung
des generalisierten Krampfanfalls. Er besteht dann, wenn der
Anfall nicht nach kurzer Zeit vergeht oder innerhalb kurzer
Abstände mehrere Anfälle auftreten (z.\,B. 2 oder mehr Anfälle
innerhalb von 24 Stunden). Der Krampf muss dann mit Medikamenten
durchbrochen werden.
}
{
\begin{itemize}
  \item Lange Dauer
  \item Wiederholtes Auftreten ($\geq$\,2\,/\,24\,h)
  \item Akute \Ca{Lebensgefahr!}
\end{itemize}
}{}{}{}

\VARIANT{VAR-RETTUNGSDIENST}
{\Fazit{Im Rettungsdienst gilt praktisch jeder Krampfanfall,
den das Personal selber miterlebt, als Status-Verdächtig.}}



\MassnahmenMK{Spezielle Maßnahmen}{Krampfender
Patient}{m:zerebraler-krampfanfall-bestendend}{\Ca{Vitale Bedrohung!}

\EE{Taktik}: \Speech{Verletzungen vermeiden,
Vitalfunktionen sichern, Abklären ob Status epilepticus
besteht, Klärung des Herganges (Fremdanamnese!),
Differentialdiagnosen abklären}
}

\MassnahmenMK{Spezielle Maßnahmen}{Nicht-(mehr) krampfender
Patient}{m:zerebraler-krampfanfall-status-post}{Neuerlicher Krampfanfall möglich!

\EE{Taktik}: \Speech{Abklären ob Status epilepticus
besteht, Klärung des Herganges (Fremdanamnese!),
Differentialdiagnosen abklären}
}
% \MassnahmenNG{Spezielle Maßnahmen}{Krampfender
% Patient}{\label{m:zerebraler-krampfanfall-bestendend}}{ NACA
% \VonBis{II}{V}:
% \Ca{Vitale Bedrohung!}
% 
% \EE{Taktik}: \Speech{Verletzungen vermeiden,
% Vitalfunktionen sichern, Abklären ob Status epilepticus
% besteht, Klärung des Herganges (Fremdanamnese!),
% Differentialdiagnosen abklären}
% 
% \begin{itemize}
%   \item \TxMassVitBes
%   \begin{itemize}
%     \item Lagerung: sichern. In der Nachschlafphase stabile
%     Seitenlage
%     \begin{itemize}
%       \item Patient vor Verletzungen schützen, Möbel etc. entfernen oder polstern
%       \item NICHT festhalten -- aber sichern, Eigenschutz
%       beachten!\footnote{\textbf{Sichern}: Patient darf
%       nicht von der Trage oder aus dem Bett fallen, der Kopf
%       darf nicht am Boden oder Tischbein anschlagen, etc.}
%     \end{itemize}
%   \end{itemize}
%   \item Besondere Diagnostik:
%   \begin{itemize}
%     \item Neurocheck inkl. Messung des \EE{Blutzuckers}!
%     \item Traumacheck
%     \item Nach \E{jedem Krampanfall}: Einschätzungsblock
%     (\FullPrettyRef{m:einschaetzungsblock}) inkl.
%     Trauma-Untersuchungen
%   \end{itemize}
%   \item Reizabgeschirmter Transport
%   \item Transportentscheidung: Abt.f. Innere Medizin. (evtl. auch
%   Intensivüberwachung)
% \end{itemize}
% }
% 
% \MassnahmenNG{Spezielle Maßnahmen}{Nicht-(mehr) krampfender
% Patient}{\label{m:zerebraler-krampfanfall-status-post}}{
% NACA \VonBis{II}{III}, neuerlicher Krampfanfall möglich!
% 
% \EE{Taktik}: \Speech{Abklären ob Status epilepticus
% besteht, Klärung des Herganges (Fremdanamnese!),
% Differentialdiagnosen abklären}
% \begin{itemize}
%   \item Umstände klären und entscheiden ob \EE{Notarzt}
%   notwendig ist:
%   \begin{itemize}
%     \item 2 oder mehr Krampfanfälle innerhalb von 24 Stunden
%     \item Anfälle \enquote{erheblich anders} als sonst üblich
%     \item Erster Krampfanfall überhaupt
%   \end{itemize}
%   \item Besondere Diagnostik:
%   \begin{itemize}
%     \item Neurocheck inkl. Messung des \EE{Blutzuckers}!
%     \item Traumacheck
%   \end{itemize}
%   \item Reizabgeschirmter Transport
%   \item Transportentscheidung: Abt.f. Innere Medizin.
% \end{itemize}
% }


\subsection{Besondere Krampfanfälle}

\Layout{60}{Fieberkrampf}
{%
\label{topic:fieberkrampf}%
\index{Fieberkrampf}%
\index{Krampfanfall!Fieberkrampf}%
%
\EE{Kinder} unter 5~Jahren können bei schnellem, hohen
Fieberanstieg ({\textgreater}\,39{\DegC}) 
\E{Fieberkrämpfe} entwickeln. Sie entsprechen einem vom
Gehirn ausgehenden Krampfanfall und sind grundsätzlich genau
so zu behandeln.

Mitunter ist die Neigung des Kindes zu Fieberkrämpfen schon
bekannt und es wurden bereits spezielle Zäpfchen für den
Bedarfsfall verschrieben (und bei Eintreffen eventuell schon
von den Eltern gegeben).
}{
\begin{itemize}
  \item Kinder {\textless} 4 Jahre
  \item Hohes Fieber
  \item Wie generalisierter Anfall
\end{itemize}
}
{./images/gabriel-sebastian-aass/fieberthermometer.jpg}
{}
{GABS}



\Layout{0}{Eklamptischer Krampfanfall während der Schwangerschaft}
{%
In Folge einer Präeklampsie (EPH-Gestose)
kann es zu einem sog. \E{eklamptischen Krampfanfall}
kommen. Dieser wird unter \FullPrettyRef{topic:eklampsie}
besprochen.
}{
\begin{itemize}
    \item Folge einer Präemklampsie.
    \item \FullPrettyRef{topic:eklampsie}.
\end{itemize}
}{}{}{}





% \clearpage % TEMP temp clearpage
\section{Lumbago, Lumbo-ischialgie und Bandscheibenvorfall} 
\HeadingSub{\Z{(Discusprolaps)}} 
\label{topic:lumbago}

\noindent Unter Ischialgie (\T{Lumbalgie}, Lumboischialgie,
\emph{\enquote{\T{Lumbago}}}) versteht man Schmerzen im
Versorgungsbereich des Ischias-Nerv, welche u. a. durch eine
Entzündung oder durch Reizung bzw. Kompression des Nervs
oder seiner Wurzeln verursacht werden können.

\Layout{0}{Ursachen}{%
können eine Reizung bzw. Druck
im Bereich der Lendenwirbel bzw. der Kreuzbeinwirbel,
ein Bandscheibenvorfall, sowie andere Erkrankungen der Wirbelsäule, aber auch
ein Unfall sein. 

\bigskip

\bigskip

}{
\begin{compactitem}
  \item Nervenreizung
  \item Bandscheibenvorfall
  \item Erkrankungen der Wirbelsäule
  \item Unfall
\end{compactitem}
}{}{}{}

\Layout{0}{\TxSymptome}{%
Das Leitsymptom sind Schmerzen in der Lendengegend, die in
das betroffene Bein bis zum Fußaußenrand ausstrahlen,
Schonhaltung, lokale Druck- und Klopfempfindlichkeit der
Wirbelsäule, Sensibilitätsstörungen und eventuell Lähmung
der Zehenmuskulatur. Je nach Schwere und zeitlicher
Entwicklung (langsames/plötzliches Auftreten) bezeichnet man
das Zustandsbild als Lumboischialgie oder Lumbago
(Hexenschuss). Gegebenenfalls ist im Krankenhaus ein
Bandscheibenvorfall abzuklären.
}{
\begin{itemize}
  \item Schmerzen in Lendengegend
  \item Evtl. Austrahlung in Bein
  \item Gefühlsstörungen
\end{itemize}
}{}{}{}


